% ══════════════════════════════════════════════════════════════════════════
%  Chapter 5: The Walsh--Hadamard Transform
% ══════════════════════════════════════════════════════════════════════════

\chapter{The Walsh--Hadamard Transform}\label{ch:wht}

This chapter develops the Walsh--Hadamard Transform (WHT) as the
computational realization of the Boolean Fourier expansion. We define
Hadamard matrices, establish their properties, describe the fast
algorithm (FWHT), and show how the WHT computes Fourier coefficients.

Historical references: \citet{walsh1923closed} for Walsh functions,
\citet{hadamard1893resolution} for Hadamard's original construction,
\citet{shanks1969computation} for the fast algorithm,
\citet{beauchamp1984walsh} and \citet{horadam2007hadamard} for
comprehensive treatments.

% ──────────────────────────────────────────────────────────────────────
\section{Hadamard Matrices}\label{sec:hadamard-def}

\begin{definition}[Hadamard matrix, Sylvester construction]\label{def:hadamard}
Define recursively:
\begin{equation}\label{eq:hadamard-recursion}
  H_0 = (1), \qquad
  H_1 = \begin{pmatrix} 1 & 1 \\ 1 & -1 \end{pmatrix}, \qquad
  H_n = H_1 \otimes H_{n-1},
\end{equation}
where $\otimes$ is the Kronecker product (\Cref{sec:kronecker}).
The matrix $H_n$ is $2^n \times 2^n$.
\end{definition}

Equivalently, $H_n = H_1^{\otimes n}$ (the $n$-fold Kronecker product of
$H_1$ with itself).

\begin{example}\label{ex:hadamard-small}
\[
H_2 = H_1 \otimes H_1 =
\begin{pmatrix}
  1 &  1 &  1 &  1 \\
  1 & -1 &  1 & -1 \\
  1 &  1 & -1 & -1 \\
  1 & -1 & -1 &  1
\end{pmatrix}.
\]
Every entry is $\pm 1$.
\end{example}


% ──────────────────────────────────────────────────────────────────────
\section{Properties of Hadamard Matrices}\label{sec:hadamard-props}

\begin{theorem}[Properties of $H_n$]\label{thm:hadamard-props}
Let $N = 2^n$. Then:
\begin{enumerate}[label=(\roman*)]
  \item \emph{All entries are $\pm 1$.}
  \item \emph{Symmetry:} $H_n = H_n^\top$.
  \item \emph{Orthogonality:} $H_n H_n^\top = N \cdot I_N$.
    Equivalently, $H_n^{-1} = \frac{1}{N} H_n$.
  \item \emph{Self-inverse up to scaling:} $H_n^2 = N \cdot I_N$.
  \item \emph{Rows are parity functions.} Row $r$ of $H_n$ (with $r$
    expressed as a subset $S \subseteq [n]$ via its binary
    representation) is the vector
    $(\chi_S(x))_{x \in \pmone^n}$ (in a specific ordering of
    the hypercube vertices).
\end{enumerate}
\end{theorem}

\begin{proof}
\emph{(i)} follows from the recursive definition: $H_1$ has $\pm 1$
entries, and the Kronecker product of $\pm 1$ matrices has $\pm 1$ entries.

\emph{(ii)} Since $H_1$ is symmetric and $(A \otimes B)^\top = A^\top \otimes B^\top$,
$H_n^\top = H_1^\top \otimes H_{n-1}^\top = H_1 \otimes H_{n-1} = H_n$
(by induction).

\emph{(iii)} By the mixed-product rule for Kronecker products:
$H_n H_n^\top = (H_1 \otimes H_{n-1})(H_1^\top \otimes H_{n-1}^\top)
= (H_1 H_1^\top) \otimes (H_{n-1} H_{n-1}^\top)$.
Since $H_1 H_1^\top = 2I_2$, by induction:
$H_n H_n^\top = 2I_2 \otimes 2^{n-1} I_{2^{n-1}} = 2^n I_{2^n} = N \cdot I_N$.

\emph{(iv)} follows from (ii) and (iii): $H_n^2 = H_n H_n^\top = N \cdot I_N$.

\emph{(v)} We establish this by identifying each row with a parity
function. Number the rows and columns from $0$ to $N-1$. For row $r$
and column $c$, write $r$ and $c$ in binary: $r = (r_{n-1}, \ldots, r_0)$
and $c = (c_{n-1}, \ldots, c_0)$. Then the $(r,c)$ entry of $H_n$ is:
\begin{equation}\label{eq:hadamard-entry}
  (H_n)_{r,c} = (-1)^{\sum_{i=0}^{n-1} r_i c_i}
  = (-1)^{r \cdot c},
\end{equation}
where $r \cdot c$ is the bitwise inner product mod~2.

Now identify the column index $c$ with the hypercube point
$x = ((-1)^{c_0}, (-1)^{c_1}, \ldots, (-1)^{c_{n-1}}) \in \pmone^n$,
and the row index $r$ with the subset $S = \{i+1 : r_i = 1\}$. Then:
\[
  (H_n)_{r,c} = (-1)^{\sum_{i \in S} c_{i-1}} = \prod_{i \in S} (-1)^{c_{i-1}} = \prod_{i \in S} x_i = \chi_S(x).
\]
So row $r$ (corresponding to subset $S$) contains the values of $\chi_S$
evaluated at all $N$ hypercube points, in the binary-counting order.
\end{proof}


% ──────────────────────────────────────────────────────────────────────
\section{The Walsh--Hadamard Transform}\label{sec:wht-def}

\begin{definition}[Walsh--Hadamard Transform]\label{def:wht}
For a vector $v \in \R^N$ (with $N = 2^n$), the \emph{Walsh--Hadamard
Transform} is
\begin{equation}\label{eq:wht}
  \hat{v} = H_n \cdot v.
\end{equation}
The \emph{normalized} WHT is $\hat{v} = \frac{1}{N} H_n \cdot v$.
The inverse transform is $v = \frac{1}{N} H_n \cdot \hat{v}$
(since $H_n^{-1} = \frac{1}{N} H_n$).
\end{definition}

\begin{theorem}[WHT computes Fourier coefficients]\label{thm:wht-fourier}
Let $f \colon \pmone^n \to \R$ and let $v \in \R^N$ be the vector of
function values in binary-counting order:
$v_c = f((-1)^{c_0}, \ldots, (-1)^{c_{n-1}})$ for
$c = 0, 1, \ldots, N-1$.

Then the normalized WHT $\hat{v} = \frac{1}{N} H_n v$ gives the Fourier
coefficients:
\[
  \hat{v}_r = \fhat(S_r),
\]
where $S_r$ is the subset corresponding to the binary representation of $r$.
\end{theorem}

\begin{proof}
\[
  \hat{v}_r = \frac{1}{N} \sum_{c=0}^{N-1} (H_n)_{r,c} \cdot v_c
  = \frac{1}{N} \sum_{x \in \pmone^n} \chi_{S_r}(x) \cdot f(x)
  = \E_x[f(x)\, \chi_{S_r}(x)] = \fhat(S_r).
\]
The second equality uses property (v) of \Cref{thm:hadamard-props}.
\end{proof}

\begin{remark}[Self-inverse property and the thesis]\label{rem:self-inverse}
Since $H_n^2 = N \cdot I$, the WHT is (up to scaling) its own inverse.
In the thesis, this means:
\begin{itemize}
  \item \emph{Forward:} Given binary weights $w \in \pmone^N$, compute
    Fourier coefficients via $\hat{w} = \frac{1}{N} H_n w$.
  \item \emph{Inverse:} Given spectral coefficients $\hat{w}$, reconstruct
    spatial weights via $w = H_n \hat{w}$ (then apply sign).
\end{itemize}
Both directions use the same matrix, just with different scaling.
\end{remark}


% ──────────────────────────────────────────────────────────────────────
\section{The Fast Walsh--Hadamard Transform}\label{sec:fwht}

The naive matrix-vector multiplication $H_n v$ costs $O(N^2)$ operations.
The \emph{Fast Walsh--Hadamard Transform} (FWHT) exploits the recursive
Kronecker structure to compute $H_n v$ in $O(N \log N) = O(N \cdot n)$
operations, using only additions and subtractions (no multiplications).

\subsection{The Butterfly Algorithm}\label{sec:butterfly}

The recursion $H_n = H_1 \otimes H_{n-1}$ means:
\[
  H_n v =
  \begin{pmatrix} I \otimes H_{n-1} \\ I \otimes H_{n-1} \end{pmatrix}
  \begin{pmatrix} v_{\text{top}} + v_{\text{bot}} \\ v_{\text{top}} - v_{\text{bot}} \end{pmatrix}
  \quad \text{(not quite right---let us be precise.)}
\]

More carefully, write $v = \binom{v_0}{v_1}$ where $v_0, v_1 \in \R^{N/2}$
are the first and second halves. Then:
\[
  H_n v = (H_1 \otimes H_{n-1}) \binom{v_0}{v_1}
  = \binom{H_{n-1}(v_0 + v_1)}{H_{n-1}(v_0 - v_1)}.
\]
Each of the two recursive calls operates on a vector of half the size.
This gives the recurrence $T(N) = 2T(N/2) + O(N)$, solving to
$T(N) = O(N \log N)$.

\begin{algorithm}[H]
\caption{Fast Walsh--Hadamard Transform (in-place)}\label{alg:fwht}
\begin{algorithmic}[1]
\Require $v \in \R^N$ with $N = 2^n$
\For{$\ell = 0, 1, \ldots, n-1$}
  \State $h \gets 2^\ell$ \Comment{half-block size}
  \For{$j = 0, 2h, 4h, \ldots, N - 2h$}
    \For{$i = j, j+1, \ldots, j + h - 1$}
      \State $a \gets v[i]$
      \State $b \gets v[i + h]$
      \State $v[i] \gets a + b$
      \State $v[i + h] \gets a - b$
    \EndFor
  \EndFor
\EndFor
\State \Return $v$ \Comment{$v$ now contains $H_n$ times the original $v$}
\end{algorithmic}
\end{algorithm}

\paragraph{Operation count.}
The triple loop performs exactly $n \cdot N/2$ additions and
$n \cdot N/2$ subtractions, for a total of $nN = N \log_2 N$ addition/subtraction
operations. \emph{No multiplications are needed.}

This is the crucial property for the thesis: when the input vector $v$
has integer entries (e.g., $v \in \pmone^N$), the entire FWHT can be
performed in integer arithmetic. The output entries are integers in
the range $[-N, N]$.

\begin{example}[FWHT of a small vector]\label{ex:fwht}
Let $n = 2$, $N = 4$, and $v = (1, -1, -1, 1)^\top$
(corresponding to the function $f(x_1, x_2) = x_1 x_2$).

\textbf{Step $\ell = 0$} ($h = 1$):
\begin{align*}
  (v[0], v[1]) &\gets (1 + (-1),\; 1 - (-1)) = (0, 2), \\
  (v[2], v[3]) &\gets (-1 + 1,\; -1 - 1) = (0, -2).
\end{align*}
After step 0: $v = (0, 2, 0, -2)$.

\textbf{Step $\ell = 1$} ($h = 2$):
\begin{align*}
  (v[0], v[2]) &\gets (0 + 0,\; 0 - 0) = (0, 0), \\
  (v[1], v[3]) &\gets (2 + (-2),\; 2 - (-2)) = (0, 4).
\end{align*}
After step 1: $v = (0, 0, 0, 4)$.

So $H_2 v = (0, 0, 0, 4)$. The normalized WHT is
$\frac{1}{4}(0, 0, 0, 4) = (0, 0, 0, 1)$.

Reading off Fourier coefficients: $\fhat(\emptyset) = 0$,
$\fhat(\{1\}) = 0$, $\fhat(\{2\}) = 0$, $\fhat(\{1,2\}) = 1$.
Indeed, $f = x_1 x_2 = \chi_{\{1,2\}}$. \checkmark
\end{example}


% ──────────────────────────────────────────────────────────────────────
\section{Comparison with the FFT}\label{sec:fwht-vs-fft}

The FWHT is the Boolean analogue of the Fast Fourier Transform (FFT).
The following table highlights the parallels:

\begin{center}
\begin{tabular}{lll}
\toprule
\textbf{Property} & \textbf{FFT (classical)} & \textbf{FWHT (Boolean)} \\
\midrule
Domain & $\Z/N\Z$ or $[0, 2\pi)$ & $\pmone^n$ \\
Basis & $e^{2\pi i k/N}$ & $\chi_S(x) = \prod_{i \in S} x_i$ \\
Basis values & Complex roots of unity & $\pm 1$ \\
Matrix & DFT matrix $F_N$ & Hadamard matrix $H_n$ \\
Self-inverse? & $F_N^{-1} = \frac{1}{N}\overline{F_N}$ & $H_n^{-1} = \frac{1}{N} H_n$ \\
Fast algorithm & $O(N \log N)$ complex & $O(N \log N)$ real \\
Arithmetic & Complex multiply-add & Integer add-subtract \\
Convolution & Circular convolution & Dyadic convolution \\
\bottomrule
\end{tabular}
\end{center}

The FWHT is simpler than the FFT: it uses only additions and
subtractions (no twiddle factors, no complex arithmetic). This makes it
particularly attractive for hardware implementation and for the
binary neural network setting of the thesis.


% ──────────────────────────────────────────────────────────────────────
\section{The Convolution Theorem via the WHT}\label{sec:wht-convolution}

The convolution theorem (\Cref{thm:convolution}) becomes a matrix
identity in the WHT:

\begin{proposition}\label{prop:wht-convolution}
Let $v, w \in \R^N$ represent functions on $\pmone^n$. Define the
\emph{dyadic convolution} vector $(v \circledast w)_c = \frac{1}{N}\sum_{d} v_d \cdot w_{c \oplus d}$,
where $\oplus$ is bitwise XOR. Then:
\[
  H_n(v \circledast w) = \frac{1}{N} (H_n v) \odot (H_n w),
\]
where $\odot$ is componentwise multiplication.
\end{proposition}

This means dyadic convolution can be performed in $O(N \log N)$ time:
transform both vectors, multiply componentwise, inverse-transform.


% ──────────────────────────────────────────────────────────────────────
\section{The WHT as a Neural Network Layer}\label{sec:wht-layer}

A key observation in the thesis is that the WHT (or its inverse) can
serve as a building block for neural network layers.

\paragraph{Hadamard layer.}
A \emph{Hadamard layer} computes $y = \frac{1}{\sqrt{N}} H_n \cdot \text{diag}(s) \cdot x$,
where $s \in \pmone^N$ is a learnable sign vector. The sign-flip
$\text{diag}(s) \cdot x$ multiplies each component of $x$ by $\pm 1$
(a componentwise sign change). The Hadamard multiplication $H_n$ then
mixes all components.

The normalization $1/\sqrt{N}$ ensures that $\norm{y}_2 = \norm{x}_2$
(since both $\text{diag}(s)$ and $\frac{1}{\sqrt{N}} H_n$ preserve the
$\ell_2$ norm---$\text{diag}(s)$ by $|s_i| = 1$, and
$\frac{1}{\sqrt{N}} H_n$ by orthogonality).

\paragraph{Spectral reconstruction.}
The thesis's \texttt{SpectralLinear} layer works differently: it stores
spectral coefficients $\hat{w}$ and reconstructs the spatial weight
vector via $w = H_n \hat{w}$ (using the FWHT), then applies $\sgn(w)$
to get binary weights. The reconstruction costs $O(N \log N)$ per row
per training step, but is not needed at inference (the binary weights are
stored directly).
